\section{Component categories and consequences}
\label{sec:consequences}

Throughout this section assume either Hypothesis~\ref{hyp:A} or
Hypothesis~\ref{hyp:B}. We apply Theorem~\ref{thm:full-support} to the spectral action.
First we identify each localized automorphic category.  We then derive the
finite-field arithmetic formula and construct Hecke eigenobjects with
coherent structure.
Throughout this section put
\[
 \mathcal C_s=\Shv_{\Nilp}(\Bun_G),\qquad
 \mathcal A=\QCoh(Y),\qquad
 \mathcal D=\IndCoh_{\Nilp}(Y),
\]
and denote the restricted Langlands equivalence by
\(\mathbb L_G^{\mathrm{restr}}:\mathcal C_s\xrightarrow{\sim}\mathcal D\).
Its \(\mathcal A\)-linear structure is the one furnished by
\cite[Lemma~1.3.7]{GR}.

\subsection{The category on a connected component}

For an open-and-closed component \(i:Z\hookrightarrow Y\), define
\[
 \mathcal C_{s,Z}:=
 \QCoh(Z)\underset{\QCoh(Y)}\otimes\mathcal C_s.
\]
Here the tensor product is the tensor product of presentable dg categories.
The idempotent quasi-coherent object \(e_Z=i_*\mathcal O_Z\) acts as the
projection to this summand.

\begin{corollary}\label{cor:components}
Under either Hypothesis~\ref{hyp:A} or Hypothesis~\ref{hyp:B}, every
connected component \(Z\subset Y\) has a nonzero localized automorphic
category.  The restriction of the given spectral-linear Langlands functor
induces an equivalence
\[
 \mathcal C_{s,Z}\xrightarrow{\sim}\IndCoh_{\Nilp}(Z).
\]
\end{corollary}
\begin{proof}
The decomposition \(Y=Z\sqcup(Y\setminus Z)\) decomposes quasi-coherent
and nilpotent ind-coherent sheaves into their two corresponding summands.
Consequently
\[
 \QCoh(Z)\underset{\mathcal A}\otimes\mathcal D
 \simeq\IndCoh_{\Nilp}(Z).
\]
Tensor the \(\mathcal A\)-linear equivalence of
Theorem~\ref{thm:full-support} with \(\QCoh(Z)\) over \(\mathcal A\).
This proves the asserted identification by the specified functor.

Every component contains a semisimple \(E\)-point
\cite[Proposition~3.7.2 and Corollary~3.7.4]{AGKRRV}.
Pullback to this point shows that \(\mathcal O_Z\ne0\).
The fully faithful embedding
\(\QCoh(Z)\hookrightarrow\IndCoh_{\Nilp}(Z)\)
therefore proves nonvanishing.
\end{proof}

The component index and the parameter itself have different roles.
Two geometric parameters lie in the same component precisely when their
semisimplifications are isomorphic
\cite[Proposition~3.7.2]{AGKRRV}.
The eigenobject construction below uses the actual parameter, including
its extension data.

\subsection{The arithmetic formula}

Assume Hypothesis~\ref{hyp:A}.  The finite-field models \(X_0\) and
\(G_0\) determine geometric \(q\)-Frobenius.  Its pullback action on
\(\QLisse(X)\) induces an automorphism \(\Frob\) of \(Y\).
Define its derived fixed-point stack by
\[
 Y^{\mathrm{arithm}}:=Y^{\Frob}
   =Y\underset{Y\times_EY}{\overset{\mathrm R}\times}Y,
\]
where the two maps to \(Y\times_EY\) are the diagonal and the graph of
\(\Frob\).  In particular, an \(E\)-point includes a Weil structure on
the geometric local system.  These are the conventions of
\cite[\S24.1]{AGKRRV} and \cite[\S1.5.1]{GR}.

\begin{corollary}\label{cor:arithmetic}
For the finite-field datum of Hypothesis~\ref{hyp:A}, there is a canonical
isomorphism
\[
 \Funct_c\bigl(\Bun_{G_0}(\mathbb F_q),E\bigr)
 \simeq
 \Gamma^{\IndCoh}\bigl(Y^{\mathrm{arithm}},
                         \omega_{Y^{\mathrm{arithm}}}\bigr).
\]
The left side is the vector space of compactly supported functions on
isomorphism classes of rational bundles, regarded as a complex concentrated
in degree zero.  The right side uses the ind-coherent dualizing object
and ind-coherent global sections on the derived fixed-point stack.
\end{corollary}
\begin{proof}
The primed union \(Y'\) is Frobenius-stable. The spectral trace
calculation of Beraldo--Lin--Reeves
\cite[\S\S6.4.12--6.4.13]{BeraldoLinReeves} enters the known primed
formula \cite[\S\S1.5.3--1.5.5 and Corollary~1.5.6]{GR}:
\[
 \Funct_c\bigl(\Bun_{G_0}(\mathbb F_q),E\bigr)
 \simeq
 \Gamma^{\IndCoh}\bigl((Y')^{\Frob},
                              \omega_{(Y')^{\Frob}}\bigr).
\]
Its automorphic trace input is
\cite[Theorem~0.2.6]{AGKRRVTrace}, with the endofunctor given by
pushforward along geometric Frobenius on \(\Bun_G\).
By Theorem~\ref{thm:full-support}, the actual inclusion \(Y'\hookrightarrow
Y\) is equality.  Taking derived fixed points identifies
\((Y')^{\Frob}\) with \(Y^{\mathrm{arithm}}\), and transports the
same dualizing object and global-sections functor.  Substitution gives
the result.
\end{proof}

Keeping \(\Gamma^{\IndCoh}\) explicit avoids a convention in
\cite[\S1.5.3]{GR}: its ordinary-global-sections notation uses the image
of the dualizing object under the dual of
\(\Upsilon:\QCoh\to\IndCoh\).
No classical truncation of \(Y^{\mathrm{arithm}}\) or replacement of its
dualizing object by its structure sheaf is made here.

\subsection{Hecke eigenobjects with coherent structure}

We first record the categorical argument that supplies the eigenobject.
It will also give its tensor compatibilities.

\begin{lemma}\label{lem:parameter-adjoint}
Let \(\mathcal A_0\) be a compactly generated, presentable, stable,
\(E\)-linear symmetric monoidal category whose compact objects are
dualizable.  Assume its tensor product preserves colimits separately.
Let \(a:\mathcal A_0\to\Vect_E\) be a continuous \(E\)-linear symmetric monoidal
functor, and give \(\Vect_E\) its \(\mathcal A_0\)-action through \(a\).
The right adjoint \(R_a\) of \(a\) is continuous and carries a coherent
\(\mathcal A_0\)-linear structure.  Moreover, \(R_a(E)\ne0\).
\end{lemma}
\begin{proof}
Presentability gives the right adjoint.  Every compact object of
\(\mathcal A_0\) is dualizable, so its image under \(a\) is a perfect
complex of \(E\)-vector spaces.  Thus \(a\) preserves compact objects.
Testing against compact generators and using adjunction shows that
\(R_a\) preserves filtered colimits.  The right adjoint is exact, since
both categories are stable.  Exactness and preservation of filtered
colimits imply preservation of all colimits.

The adjunction has a canonical projection morphism
\[
 b\otimes R_a(V)\longrightarrow R_a\bigl(a(b)\otimes V\bigr).
\]
For dualizable \(b\), its invertibility follows by testing maps out of
an arbitrary \(c\in\mathcal A_0\):
\begin{align*}
 \RHom_{\mathcal A_0}(c,b\otimes R_a(V))
 &\simeq\RHom_{\mathcal A_0}(b^\vee\otimes c,R_a(V))\\
 &\simeq\RHom_E(a(b)^\vee\otimes a(c),V)\\
 &\simeq\RHom_E(a(c),a(b)\otimes V).
\end{align*}
Both sides of the projection morphism preserve colimits in \(b\).
Compact generation therefore proves invertibility for every \(b\).
These morphisms are the adjunction mates of the module structure on
\(a\); hence their unit and associativity compatibilities hold before
invertibility is imposed.  They now give a coherent module-functor
structure on \(R_a\).
Finally,
\[
 \RHom_{\mathcal A_0}(\mathbf1,R_a(E))
 \simeq\RHom_E(E,E)\simeq E,
\]
so \(R_a(E)\ne0\).  This last argument does not require a compact unit.
\end{proof}

For our spectral stack, \(\mathcal A=\QCoh(Y)\) satisfies the hypotheses
of Lemma~\ref{lem:parameter-adjoint} by
\cite[\S7.9 and Corollary~C.4.4]{AGKRRV}.
On the disjoint union of components, the compact generators have finite
component support.  We use the continuous, fully faithful \(\mathcal A\)-linear functor
\[
 \jmath_Y:\QCoh(Y)
 \xrightarrow[\sim]{\Upsilon_Y}\IndCoh_{\{0\}}(Y)
 \hookrightarrow\IndCoh_{\Nilp}(Y).
\]
Here \(\Upsilon_Y(Q)=Q\otimes\omega_Y\), as in
\cite[\S1.3.1, Equation~(1.2)]{GR}.
Its essential image has zero singular support and therefore lies in the
nilpotent category.

For a finite set \(I\), put \(\mathcal Q_I=\QLisse(X)^{\otimes I}\).
Given \(V\in\Rep(\check G)^{\otimes I}\), let \(H_I(V,-)\) denote the
multi-leg Hecke functor and let
\(\mathcal E_V^I\in\mathcal A\otimes\mathcal Q_I\)
be the universal evaluation object.  For a parameter
\(\sigma:\Spec(E)\to Y\), define its evaluation local system by
\[
 \sigma^I(V):=(\sigma^*\otimes\id)(\mathcal E_V^I)
 \in\mathcal Q_I.
\]
For a family of representations \((V_i)_{i\in I}\), this is the exterior
product of the local systems \(\sigma(V_i)\) on the corresponding factors
of \(X^I\).

\begin{definition}\label{def:eigenstructure}
A Hecke eigenobject of eigenvalue \(\sigma\) with coherent structure is an
object \(M\in\mathcal C_s\) equipped, for all finite sets \(I\), with
isomorphisms
\[
 H_I(V,M)\simeq M\boxtimes\sigma^I(V)
\]
natural in \(V\in\Rep(\check G)^{\otimes I}\) and compatible with units,
tensor products, permutations of the legs and collisions of legs, with
all compatibilities imposed homotopy-coherently
\cite[\S15.2]{AGKRRV}.
\end{definition}

\begin{corollary}\label{cor:eigenobjects}
Under either Hypothesis~\ref{hyp:A} or Hypothesis~\ref{hyp:B}, every
parameter \(\sigma:\Spec(E)\to Y\) admits a nonzero object
\(M_\sigma\in\Shv_{\Nilp}(\Bun_G)\) with a coherent Hecke eigenstructure
of eigenvalue \(\sigma\).  If \(\sigma\) belongs to a component \(Z\),
then \(M_\sigma\in\mathcal C_{s,Z}\).
The assertion concerns the ind-constructible dg category and includes
no boundedness, perversity or Weil structure.
\end{corollary}
\begin{proof}
Apply Lemma~\ref{lem:parameter-adjoint} to
\(a=\sigma^*:\mathcal A\to\Vect_E\), and write \(\sigma_*\) for its
continuous right adjoint.  The composite
\[
 \Phi_\sigma:
 \Vect_\sigma\xrightarrow{\sigma_*}\mathcal A
 \xrightarrow{\jmath_Y}\mathcal D
 \xrightarrow{(\mathbb L_G^{\mathrm{restr}})^{-1}}\mathcal C_s
\]
is a continuous \(\mathcal A\)-linear functor.  Here \(\Vect_\sigma\)
denotes \(\Vect_E\) with its action through \(\sigma^*\).
Set \(M_\sigma=\Phi_\sigma(E)\).
It is nonzero by the lemma, full faithfulness of
\(\jmath_Y\), and the Langlands equivalence.
If \(\sigma\) lies in \(Z\), the idempotent \(e_Z\) pulls back to \(E\),
while its complementary idempotent pulls back to zero.
Linearity therefore places \(M_\sigma\) in \(\mathcal C_{s,Z}\).

To verify the eigenstructure, fix a finite set \(I\).
The action of \(\mathcal E_V^I\) on the automorphic category is the
multi-leg Hecke functor \(H_I(V,-)\)
\cite[\S\S14.2--14.3]{AGKRRV}.
Extend \(\Phi_\sigma\) by the identity on \(\mathcal Q_I\).
Its module structure gives
\[
 H_I(V,M_\sigma)
 \simeq(\Phi_\sigma\otimes\id)\bigl(\sigma^I(V)\bigr)
 \simeq M_\sigma\boxtimes\sigma^I(V).
\]
The second equivalence uses that a continuous \(E\)-linear functor from
\(\Vect_E\) is tensoring with its value on \(E\).

These equivalences are natural in \(V\).  They preserve tensor products
and the unit by the coherent module structure of \(\Phi_\sigma\), and
are compatible with collisions of legs and maps between finite sets by
the universal symmetric monoidal evaluation system.
Thus they supply the homotopy-coherent Hecke eigenstructure of
\cite[\S15.2]{AGKRRV}.  Equivalently, restricting \(\Phi_\sigma\) along
the symmetric monoidal Ran-to-spectral action gives the eigenobject in
\cite[\S15.1]{AGKRRV}.
The evaluation uses \(\sigma\) itself, so its eigenvalue is the actual
parameter, whether or not it is semisimple.
\end{proof}
